% ==============================================================================
% Section 3: Methods
% ==============================================================================
\section{Methods}
\label{sec:methods}

\subsection{Spatial Weights}
\label{sec:weights}

We construct a Queen contiguity spatial weights matrix for all \nmunis{} municipalities using the \texttt{libpysal} library \citep{rey_pysal_2010}.
To address sub-metre geometry precision gaps in the source geoparquet that would otherwise produce false island municipalities, we employ fuzzy contiguity with a 50-metre buffer (\texttt{fuzzy\_contiguity(buffering=True, buffer=50)}) rather than exact topological intersection.
The matrix is row-standardized.
This approach yields \nislands{} island municipalities---all boundary gaps are resolved by the buffer tolerance.

\subsection{Global Spatial Autocorrelation}
\label{sec:global}

Global Moran's $I$ quantifies the overall degree of spatial clustering:

\begin{equation}
  I = \frac{n}{\sum_i \sum_j w_{ij}}
      \cdot
      \frac{\sum_i \sum_j w_{ij}(x_i - \bar{x})(x_j - \bar{x})}
           {\sum_i (x_i - \bar{x})^2}
  \label{eq:moran}
\end{equation}

\noindent where $x_i$ is the case count in municipality $i$, $w_{ij}$ are row-standardized Queen contiguity weights, and $n = \text{\nmunis{}}$.
Statistical significance is assessed via 999 conditional permutations at $\alpha = 0.05$.

\subsection{Local Spatial Autocorrelation (LISA)}
\label{sec:lisa}

Local Moran's $I_i$ decomposes the global statistic into municipality-level contributions \citep{anselin_1995}:

\begin{equation}
  I_i = \frac{(x_i - \bar{x})}{m_2} \sum_j w_{ij}(x_j - \bar{x}), \quad
  m_2 = \frac{1}{n}\sum_i (x_i - \bar{x})^2
  \label{eq:lisa}
\end{equation}

\noindent Municipalities are classified into five types based on the sign of their local statistic and the significance of their pseudo-$p$-value at $\alpha = 0.05$: High-High (HH), Low-Low (LL), High-Low (HL), Low-High (LH), and Not Significant (NS).
HH municipalities have high counts surrounded by high-count neighbors; LL municipalities have low counts surrounded by low-count neighbors.

We extract connected-component HH clusters via breadth-first search on the Queen adjacency graph, retaining only clusters of three or more contiguous HH municipalities.
LL clusters are retained in tables for completeness but are not interpreted as a substantive finding: the extreme zero-inflation of the count data (77--98\% of municipality-months record zero cases, depending on status) means that LL classification largely identifies municipalities where nothing happened surrounded by municipalities where nothing happened.

The total computational scope is \nLISAruns{} LISA runs: 4 statuses $\times$ 3 sex categories (total, male, female) $\times$ \nmonths{} months, each computed over $N = \text{\nmunis{}}$ municipalities with 999 permutations.

\subsection{Geographic Concentration}
\label{sec:concentration}

We quantify geographic concentration using the Gini coefficient:

\begin{equation}
  G = \frac{2\sum_{i=1}^{n} i \cdot x_{(i)}}{n \sum_{i=1}^{n} x_{(i)}}
      - \frac{n+1}{n}
  \label{eq:gini}
\end{equation}

\noindent and the Herfindahl-Hirschman Index (HHI):

\begin{equation}
  \text{HHI} = \sum_{i=1}^{n} s_i^2, \quad s_i = \frac{x_i}{\sum_j x_j}
  \label{eq:hhi}
\end{equation}

\noindent where $x_{(i)}$ is the $i$-th order statistic of the count distribution.
Both metrics are computed for each of the \nmonths{} monthly cross-sections per status.
We benchmark our concentration estimates against \citet{weisburd_2015}'s ``law of crime concentration,'' which posits that approximately 50\% of crime concentrates in 4--6\% of places.

\subsection{Trend Analysis}
\label{sec:trend}

We fit OLS regressions of the form $y_t = \alpha + \beta t + \varepsilon_t$ separately for each (region $\times$ status) combination, where $y_t$ is the monthly count of HH municipalities and $t = 1, \ldots, \text{\nmonths{}}$.
Slopes are expressed as municipalities per year ($\hat\beta \times 12$).
Because monthly HH counts exhibit serial correlation, standard errors are computed using the \citet{newey_west_1987} heteroskedasticity and autocorrelation consistent (HAC) estimator with a bandwidth of 12 months.

\subsection{Sex Disaggregation}
\label{sec:sex_methods}

LISA is repeated separately for male and female case counts across all statuses and months.
We compare male and female HH geographies using Jaccard similarity (intersection over union of HH municipality sets) for each status at the \referencemonth{} cross-section.
A critical limitation applies to female Located Dead: with a mean of 0.0065 cases per municipality-month and 99.4\% zero cells, the LISA algorithm lacks sufficient spatial variation to detect meaningful clusters.
Female Located Dead results are reported for completeness but should be interpreted with extreme caution.

\subsection{CED Alignment Analysis}
\label{sec:ced_methods}

As a post-hoc exercise, we extract HH dynamics for eight states explicitly named in the Committee on Enforced Disappearances' Article~34 decision \citep{CED_MEX_art34_2026}: Jalisco, Guanajuato, Coahuila, Veracruz, Nayarit, Tabasco, Nuevo Le\'{o}n, and Estado de M\'{e}xico.
For each state, we count the number of municipalities classified as HH in each month, overlay the CED's temporal claims, and compute OLS trend slopes.
No additional LISA runs are required; this analysis operates entirely on the existing \nLISAruns{} runs.
